\input{estilo}
% El archivo lógico partido en chunks de 64 MB, numerados desde 0 como el
% chunk index del texto, y repartido entre chunkservers: cada chunk está en
% tres. Los cuatro chunkservers y el reparto no están en la pizarra ni en las
% notas, que solo dicen "replicados en chunkservers". El chunk 2 va en
% verde en el archivo y en sus tres copias, para seguirlo.
\begin{document}
\begin{figura}\begin{tikzpicture}
  \def\w{1.5}\def\h{0.62}
  \fill[relleno] (0,0) rectangle (\w,4*\h);
  \fill[segundo!25] (0,\h) rectangle (\w,2*\h);
  \draw (0,0) rectangle (\w,4*\h);
  \foreach \k in {1,2,3} \draw (0,\k*\h) -- (\w,\k*\h);
  \foreach \k in {0,...,3} \node[etiq] at (\w/2,{(3.5-\k)*\h}) {chunk \k};
  \node[cod, below=3pt] at (\w/2,0) {/dir/file.txt};
  \draw[llave=4pt, draw=tinta] (-0.15,0) -- (-0.15,\h) node[midway, left=6pt, etiq] {64 MB};
  % los chunkservers, cada uno con tres chunks; el 2, en verde, en sus tres copias
  \foreach \s/\ks [count=\j] in {1/{0,1,2}, 2/{1,2,3}, 3/{2,3,0}, 4/{3,0,1}} {
    \begin{scope}[shift={(3.0,{(4.5-\j)*0.95-0.6})}]
      \node[caja, minimum width=26mm, minimum height=7.5mm] at (0.75,0) {};
      \node[etiq, font=\footnotesize, anchor=east] at (-0.6,0) {CS\textsubscript{\s}};
      \foreach \k [count=\i] in \ks {
        \ifnum\k=2 \def\f{segundo!25}\else\def\f{white}\fi
        \node[draw=tinta, fill=\f, minimum width=6mm, minimum height=4.5mm, inner sep=0pt, font=\scriptsize]
          at ({(\i-1)*0.75},0) {\k};
      }
    \end{scope}
  }
\end{tikzpicture}\end{figura}
\end{document}
